\chapter{Street's 2004 Arabic logic: traditions and syllogistic disputes}
\label{sec:p13-street-arabic-logic-2004}

Tony Street's 2004 chapter reconstructs Peripatetic logical writing
in Arabic between 750 and 1350, particularly the syllogistic debates
which issued in later teaching texts. Its central question is which
earlier systems and problems a logician addresses, not whether all
Arabic authors exemplify one doctrine. This source-bound treatment
uses the actual published chapter, pp.~523--596, including its
three appendices \cite{Street2004ArabicLogicHHL}.

\section{Historical method and a translated logical corpus}
\label{sec:p13-street-2004-method-corpus}

Street rejects two shortcuts. Madkour's assumption of agreement
on major logical questions makes a single paradigm author appear
sufficient. Rescher instead assigns logical allegiance through
general philosophical schools. Street uses the second approach
as a starting heuristic, but tests it against actual logical
references, constructions and disagreements. Pedagogical lineage
need not constitute a uniform doctrine or fixed curriculum.
Century-by-century divisions need not locate the significant
episodes of logical activity. His historical reconstruction
remains tentative, especially where manuscripts are unedited.
Ending at 1350 is not evidence that original work then ceased.
\cite[pp.~523--529]{Street2004ArabicLogicHHL}

The translation account distinguishes Syriac teaching, early
Arabic work, the Kindi and Hunayn circles, and later critical
revision. Sergius of Reshaina, Severus Sebokht and Athanasius
of Balad belong to the earlier Syriac story. Ibn al-Muqaffa's
epitome and the caliphal commission of the \emph{Topics} belong
to the early Arabic story. Patronage, contract work, religious
communities and philosophical priorities varied; this was not
one linear sequence in which every later version replaced
every earlier version. Galenic logical material was also
translated even though al-Farabi and Avicenna later acknowledged
little debt to Galen.
\cite[pp.~529--532]{Street2004ArabicLogicHHL}

Ibn Suwar's annotated \emph{Organon}, represented by Paris
arabe~2346, connects translation, textual comparison and
philosophical interpretation. Ibn al-Nadim's \emph{Fihrist},
completed in 987, supplies a bookseller's evidence about the
translation enterprise and its commentators. For the \emph{Prior
Analytics}, Street distinguishes the uncertain identity and date
of Tadari, correction by Hunayn, the Arabic/Syriac direction
of particular revisions, and the mass of sometimes incompatible
interpretative material. Neither a familiar name nor a later
edition removes those uncertainties.
\cite[pp.~532--533]{Street2004ArabicLogicHHL}

\section{Logic, meanings and the organization of knowledge}
\label{sec:p13-street-2004-language-knowledge}

Al-Farabi's account compares logic's regulation of intelligibles
with grammar's regulation of a particular language. This is
both philosophical and apologetic in the Baghdad setting,
where Abu Bishr Matta's encounter with al-Sirafi made the
utility of logic disputed. Al-Farabi also develops a context
theory: general logical preparation supports demonstrative,
dialectical, sophistical, rhetorical and poetic discourse,
distinguished by their matter, epistemic grade and aims.
The relation to psychology and knowledge is substantive;
one valid inferential shape does not make every kind of
premise demonstrative.
\cite[pp.~536--539]{Street2004ArabicLogicHHL}

Avicenna accepts logic's normative role but denies that
expressions signifying meanings define its subject matter.
His second-intention account concerns meanings classified
as subjects, predicates, universals and other logical roles
in obtaining the unknown from the known. Logic can therefore
be both an instrument and a science with a subject of its
own. Language is an accidental route for human thought,
not a requirement inherent in every intellectual operation.
His rejection of the \emph{Categories} as necessary preparation
for syllogisms and his modified placement of syllogistic
matter change the architecture of the logical treatise.
\cite[pp.~539--541]{Street2004ArabicLogicHHL}

Conception and assent are different targets of investigation.
A concept such as house does not itself affirm that houses
have a particular property. An explanatory phrase, including
definition or description, seeks the former target; proof,
including syllogism and induction, seeks the latter. Street
traces the distinction from al-Farabi's demonstration theory
to Avicenna's use of it at the beginning of logical writing.
The difference concerns knowledge and method, not simply
whether an expression is long or short.
\cite[pp.~540--542]{Street2004ArabicLogicHHL}

\section{Hypothetical inference and reduction}
\label{sec:p13-street-2004-hypothetical}

Street uses hypothetical syllogistic as a broad comparison
term, not a technical expression with identical extension
in al-Farabi and Avicenna. Al-Farabi's surviving treatment
offers the five indemonstrable schemata and does not use
them to analyze reduction to a contradiction. Avicenna
belongs to a tradition postulating a lost Aristotelian
work and recalls a longer book available in his homeland.
These are historical claims of the authors, not evidence
that a lost work has now been recovered. Street's argument
against assigning an extensive lost hypothetical treatment
to al-Farabi uses Avicenna's uncertain attribution and
Averroes' failure to invoke such a work; he does not
claim to have read the lost commentary.
\cite[pp.~543--545]{Street2004ArabicLogicHHL}

Avicenna's conditionals admit quantity and quality:
always, sometimes, never, and sometimes not. His
conjunctive syllogisms do not state the conclusion or
its contradictory explicitly in a premise; exceptive
syllogisms do. These distinctions are not identical
to the categorical/hypothetical division. Street
supplies a wholly hypothetical first-figure chain,
and a mixed case in which a conditional consequent
combines with a categorical universal to produce
a new conditional consequent.
\cite[pp.~545--547]{Street2004ArabicLogicHHL}

\begin{hbexample}{def:p13-street-2004-conditional-chain}{The shared consequent is an inferential connection}
Suppose it is always the case that if the first
specified predication holds, a second holds, and
always that if the second holds, a third holds.
The first-to-third conditional is obtained through
the shared intermediate predication. In the mixed
case, if an antecedent holds then every member of
one class belongs to another, and every member of
that second class belongs to a third; under the
same antecedent the first class belongs to the third.
The argument keeps the condition, class senses
and relevant occasions fixed. It is the stated
composition, not an identification of quantified
historical conditionals with every modern material
conditional in an unspecified temporal model.
\end{hbexample}

Street's worked reduction is Baroco: all members
of one class belong to a second, some members of
a third do not belong to the second, therefore
some of the third do not belong to the first.
The source separates a conjunctive step constructing
a conditional from the denial of that consequent
in an exceptive step. The following result checks
that structure under explicitly classical, fixed
interpretations; it does not supply Avicenna's full
modal semantics.
\cite[pp.~546--547]{Street2004ArabicLogicHHL}

\begin{hbtheorem}{thm:p13-street-2004-baroco}{proved=annex:p13-street2004-baroco}{A fixed-interpretation check of the two-stage reduction}
For unary predicates \(B,D,J\) in ordinary classical
first-order semantics, premises
\[
 \forall x(B(x)\to D(x)),\qquad
 \exists x(J(x)\land\neg D(x))
\]
entail \(\exists x(J(x)\land\neg B(x))\).
If the desired conclusion is denied, the first
premise produces \(\forall x(J(x)\to D(x))\);
the second premise contradicts that consequent.
No premise asserting a nonempty \(B\)-class is needed.
\end{hbtheorem}

Kashi's later alternative uses an exceptive with
a compound antecedent: if Zayd writes and every
writer moves fingers, Zayd moves fingers; Zayd
does not move fingers. Inferring that Zayd does
not write also requires retaining the granted
universal. Denial of the conjunction alone would
not select which conjunct fails. Street records
Kashi's account and Tusi's apparent adoption,
while reserving judgment on the depth of the
argument; it is not the Averroist critique merely
under another name.
\cite[pp.~576--577]{Street2004ArabicLogicHHL}

\section{Absolute propositions and competing conversion rules}
\label{sec:p13-street-2004-absolute-modal}

The same translated term can support different
truth conditions. Al-Farabi's absolute or existential
proposition belongs to an assertoric foundation.
Avicenna's absolute instead carries an at-least-once
temporality and is contradicted by a perpetual.
He stipulates this temporal reading for negatives
as well as affirmatives despite an ordinary negative
often suggesting never. A fixed-time negative
converts, but restricting the subject to that time
does not answer a question about all individuals
eligible for logical treatment.
\cite[pp.~547--550]{Street2004ArabicLogicHHL}

The no-human-laughing example makes the contrast
concrete. An absence of laughter at an occasion
does not imply that nothing laughing at any
other occasion is human. Alexander's and
al-Farabi's reduction through Ferio presumes
one interpretation of the negative and its
contradictory; Avicenna objects to its compatibility
with other stipulated doctrines. A modern
pointwise conversion is not a proof of conversion
for a historical proposition with different
temporal and subject conditions.
\cite[pp.~548--549]{Street2004ArabicLogicHHL}

Al-Farabi's ampliated subject and reduplicative
condition enter the dispute about a contingent
proposition's converse. Avicenna objects that
adding a condition changes the original sleeping
proposition, and that a stratagem used in
conversion must also be respected in the mixed
Barbara argument. Averroes defends an implicit
condition and distinguishes essential from
accidental modality, but Street judges the
conversion usage equivocal. This is a criticism
of the reported solutions, not a claim that
one modern necessity operator validates all of them.
\cite[pp.~549--552, 563--564]{Street2004ArabicLogicHHL}

Avicenna's six conditions distinguish unqualified
necessity, the substance's existence, its qualifying
description, a predicate condition, a definite
time and an indefinite time. The substance and
description conditions dominate his treatment.
A bachelor can cease to be unmarried as an
individual while remaining unmarried whenever
described as a bachelor; a sleeper can wake,
but not while retaining the sleeping description.
Later writers investigate definite and indefinite
temporal modes much more extensively.
\cite[pp.~550--552]{Street2004ArabicLogicHHL}

The later figure-world case separates actual
extension from possible meanings. In a world
containing only triangular figures, every actual
figure can be triangular without the meaning
of figure implying triangle. Katibi's choice
of subject interpretation helps explain rejection
of possible minors that Avicenna accepted. The
later descriptional perpetual also acquires a
descriptional at-least-once contradictory rather
than a second perpetual; naming the new mode
\emph{hiniyya} is part of that development,
not proof that the old pair was already sound.
\cite[pp.~574--576]{Street2004ArabicLogicHHL}

\section{Religious reception and the alleged clash of schools}
\label{sec:p13-street-2004-reception-schools}

Street distinguishes grammatical, theological
and juridical objections. Accounts of the
Matta--Sirafi debate disagree about its point
and setting. Similarity between theological
argument vocabulary and Stoic material is
suggestive but does not recover the route
of transmission. Ibn Khaldun's account of
reversible arguments explains a theological
worry: treating failure of an argument as
failure of the article it supports makes
criticism of premises threatening. Rejecting
that inference allows replacement of arguments
without abandonment of the articles. Street
retains the limits of this historical explanation.
\cite[pp.~554--558]{Street2004ArabicLogicHHL}

Al-Ghazali's principal contribution in this
account is practical reception: domesticating
terminology, distinguishing logical from
grammatical usage and explaining the usefulness
of logical form for juristic reasoning. His
formal treatment remains elementary compared
with the debates over Avicenna's modals.
Ibn al-Salah and Ibn Taymiyya represent
continued opposition; questioning utility,
certainty or doctrinal association does not
amount to refuting every valid syllogistic form.
\cite[pp.~558--560]{Street2004ArabicLogicHHL}

The Andalusian account includes Ibn Abdun's
reported pedagogical connection, Ibn Hazm's
apologetic project, Abu al-Salt's summary,
Ibn Bajja's Farabian commentaries and Averroes'
changing solutions. Averroes' hypothetical
critique disputes Avicenna's division and
reduction claims while defending the relevance
of Aristotle's limited treatment. Gersonides'
critical reception is distinct from Maimonides'
short logical tract. Ibn Malih's rare Arabic
response also shows that rejecting Avicenna
need not entail accepting every Averroist repair.
\cite[pp.~561--567]{Street2004ArabicLogicHHL}

Street rejects a simple Western/Eastern
opposition centred on Abu al-Barakat, Razi
and Tusi. Abu al-Barakat accepts important
Avicennan modal conversions while altering
the absolute negative; no continuing school
is thereby established. Abd al-Latif's
Farabian revival has a different, incompletely
recovered textual basis. Razi is a source
for Tusi's understanding of al-Farabi and
Sawi, not only a target. Parallel passages
and acknowledged dependence matter more
than importing a metaphysical rivalry into
logic. Tahtani's arbitration is not assumed
to synthesize two fundamentally opposed
logical systems merely from its title.
\cite[pp.~568--574]{Street2004ArabicLogicHHL}

The later handbooks emerge from continuing
revision, not unmodified adherence. Khunaji's
\emph{Jumal}, Abhari's introductory work and
Katibi's \emph{Shamsiyya} differ in depth
and emphasis. Street's criterion of Avicennan
affiliation concerns the absolute's truth
conditions, the available readings and the
conjunctive/exceptive division; post-Avicennan
is only a chronological description.
The chapter's judgments and open problems
are those of its 2004 investigation, not
a census of subsequent scholarship.
\cite[pp.~579--585]{Street2004ArabicLogicHHL}

\section{The source's modal tables and their domains}
\label{sec:p13-street-2004-modal-tables}

Appendix A uses a mode followed by predicate,
subject and a superscript quantity/quality:
\(a,e,i,o\) mean universal affirmative,
universal negative, particular affirmative
and particular negative. \(X\) is absolute,
\(A\) perpetual, \(M\) possible and \(L\)
necessary; \(w\) marks a descriptional
reading. Mode triples give major, minor,
conclusion in that order. The local sign
\(\not\equiv\) records the source's
contradictory pairing, not mere nonequivalence.
\cite[pp.~590--592]{Street2004ArabicLogicHHL}

For \(X_1\) and \(A\), the pairings are
\[
 X_1bj^a\not\equiv Abj^o,\quad
 X_1bj^e\not\equiv Abj^i,\quad
 X_1bj^i\not\equiv Abj^e,\quad
 X_1bj^o\not\equiv Abj^a.
\]
For \(X_2\), the respective contradictories
are \(Abj^i\lor Abj^o\) for \(a,e\),
and \(Abj^a\lor Abj^e\) for \(i,o\).
The source reports corresponding isomorphic
\(M_1,L\) and \(M_2,L\) tables.
These are historical statement tables with
their original \emph{Pointers} referrals,
not an independently established adequacy theorem.

The first-figure perfect mixes listed are
\(XXX,XLX,LXL,LLL,MMM,MXM,MLM\);
some second-figure proofs use \(AXA,ALA\).
Further first-figure mixes are \(LML,XMM\).
Second-figure results are \(LLL,XLL,LXL,MLL,LML\),
with \(XX,MM,XM,MX\) nonproductive;
third-figure results are
\(XXX,LLL,LXL,XLX,MMM,XMM,MXM,LML,MLM\).
Absolute negative conversion fails;
\(X_1\) affirmative universal/particular
and \(X_2\) affirmative universal/particular
convert to \(X_1\) particulars. Substituting
\(M_1,M_2\) gives the source's possible
conversion list. Necessary negatives convert
to necessary negatives; necessary affirmatives
convert to possible particulars. For purely
descriptional premises Appendix A reports
fourteen productive moods and same-kind
negative conversion and contradiction;
it then considers the mixed Barbara
\(L^wLL\). Its explanation keeps the
minor's continuing predication and the
major's description-bound duration distinct.
These additional statements retain the
source's cited scope, not a newly proved
modal adequacy claim.
\cite[pp.~590--591]{Street2004ArabicLogicHHL}

The later table instead excludes possible
minors and possible conversion; necessary
negatives convert to perpetual negatives,
and necessary affirmatives to descriptional
perpetual particulars. Descriptional perpetual
and absolute modes form the four oppositional
pairings; negative perpetual conversion
preserves that mode, while affirmative
perpetuals convert to descriptional absolute
particulars. The first-figure productivity
rule requires one of the source's seventeen
actual modes as minor; ordinarily
the major supplies the conclusion's mode,
with separate provisions for descriptional
major modes and their restrictions. Street's
printed rule calls its exceptional group
four while displaying three modes; the
missing case is not silently supplied here.
\cite[p.~592]{Street2004ArabicLogicHHL}

\subsection{The 26 later-mode renditions}

The following guide preserves all 26 items in
Appendix B without treating them as Avicenna's
own semantics. Here \(R_t(Q_x)\) is realization
of the predication \(Q\) of an individual \(x\)
at time \(t\). The unbound temporal operators
\(\exists Q_x\) and \(\forall Q_x\) abbreviate
\((\exists t)R_t(Q_x)\) and \((\forall t)R_t(Q_x)\);
they are not quantifiers over individuals.
\(TQ_x=R_T(Q_x)\) and \(SQ_x=R_S(Q_x)\)
use the source's specified time constants.
Modal prefixes are preserved: for example,
\(\exists\Box Q_x=(\exists t)\Box R_t(Q_x)\),
\(\forall\Diamond Q_x=(\forall t)\Diamond R_t(Q_x)\),
and \(\Box TQ_x=\Box R_T(Q_x)\).
The modal operator's position is therefore
not silently interchanged with a temporal quantifier.
\cite[pp.~593--595]{Street2004ArabicLogicHHL}

Each listed affirmative has the printed common
form \(\forall x[\exists A_x\supset F_x]\),
where \(A,B\) are predicates and \(\supset\)
is the rendition's material conditional.
The source suppresses an existence condition;
its restored example for necessary predication is
\[
 \forall x[(\exists t)R_t(A_x)\supset
          (\forall t)\Box R_t(A_x\supset B_x)].
\]
That declared suppression and the assumed
quantificational/modal machinery are part
of the rendition's scope, not a proof of
its historical adequacy.

\begin{longtable}{@{}r p{0.43\textwidth} p{0.43\textwidth}@{}}
\caption{The 26 later-mode renditions}\\
\toprule
Item & Source mode name & Printed consequent \(F_x\) \\
\midrule
\endfirsthead
\caption[]{The 26 later-mode renditions (continued)}\\
\toprule
Item & Source mode name & Printed consequent \(F_x\) \\
\midrule
\endhead
\midrule\multicolumn{3}{r}{\emph{continued on next page}}\\
\endfoot
\endlastfoot
1 & Absolute necessary, \(L\) & \(\forall\Box B_x\) \\
2 & General conditional, \(L^w\) & \(\forall\Box(A_x\supset B_x)\) \\
3 & Absolute temporal & \(\Box TB_x\) \\
4 & Absolute spread & \(\Box SB_x\) \\
5 & Absolute perpetual, \(A\) & \(\forall B_x\) \\
6 & General conventional, \(A^w\) & \(\forall(A_x\supset B_x)\) \\
7 & Absolute continuing, \(X^w\) & \(\exists(A_x\supset B_x)\) \\
8 & Temporal absolute & \(TB_x\) \\
9 & Spread absolute & \(SB_x\) \\
10 & General absolute, \(X_1\) & \(\exists B_x\) \\
11 & Possible continuing & \(\exists\Diamond(A_x\supset B_x)\) \\
12 & Temporal possible & \(\Diamond TB_x\) \\
13 & General possible, \(M_1\) & \(\exists\Diamond B_x\) \\
14 & Captioned perpetual possible & \(\Diamond SB_x\) \\
15 & Non-perpetual necessary & \(\forall\Box B_x\land\neg\forall B_x\) \\
16 & Special conditional & \(\forall\Box(A_x\supset B_x)\land\neg\forall B_x\) \\
17 & Temporal & \(\Box TB_x\land\neg\forall B_x\) \\
18 & Spread & \(\Box SB_x\land\neg\forall B_x\) \\
19 & Captioned non-perpetual perpetual & \(\forall B_x\land\neg\forall B_x\) \\
20 & Special conventional & \(\forall(A_x\supset B_x)\land\neg\forall B_x\) \\
21 & Non-perpetual continuing absolute & \(\exists(A_x\land B_x)\land\neg\forall B_x\) \\
22 & Non-perpetual temporal absolute & \(TB_x\land\neg\forall B_x\) \\
23 & Non-perpetual spread absolute & \(SB_x\land\neg\forall B_x\) \\
24 & Non-perpetual existential, \(X_2\) & \(\exists B_x\land\neg\forall B_x\) \\
25 & Non-necessary existential & \(\exists B_x\land\neg\forall\Box B_x\) \\
26 & Special possible, \(M_2\) & \(\exists\Diamond B_x\land\neg\forall\Box B_x\) \\
\bottomrule
\end{longtable}

The original illustrations distinguish a
substance-linked rationality claim, finger
movement during writing, an eclipse at a
specified alignment, and breathing at spread
times. They illustrate stipulated logical
readings; their wording does not establish
new biological or physical laws. In particular,
\(\forall\Box B_x\) and \(\exists\Box B_x\)
express different temporal strengths even
before comparison with the historical propositions.

Two printed items are diagnostics, not
adopted definitions of a correct system.
Item~14 is captioned perpetual and glossed
with all times, but its formula uses the
spread constant \(S\), not temporal \(\forall\).
Item~19 has \(F_x\) conjoined with its own
negation: under the stated classical reading
no individual satisfying \(\exists A_x\)
can satisfy that consequent. It therefore
cannot serve as a nonempty affirmative
mode in that literal form. The intended
repairs require the underlying source;
this guide does not invent them. Likewise,
an existentially quantified material
conditional such as item~7 can be true
at a time when its antecedent fails; its
English gloss alone does not remove that
vacuity. These qualifications prevent a
literal source rendition becoming an
unjustified historical semantics.

\begin{proofboundary}{def:p13-street-2004-formal-boundary}{Source tables are not a completed historical calculus}
The exact source has now been recovered,
including its 26 later-mode renditions.
Those renditions explicitly suppress subject
existence conditions, assume an unspecified
usual quantificational/modal machinery and
are expressly denied as an interpretation
of Avicenna's own propositions. They cannot
be substituted for that historical semantics.
The printed four/three discrepancy, the
contradictory literal rendition at item~19,
and the temporal-caption mismatch at item~14
require source reconciliation. Corrected
historically qualified renditions and justified proof
routes for the modal lists remain separate
formal work; no theorem about their whole
consistency, equivalence or completeness
is asserted by the historical narrative above.
\cite[pp.~592--596]{Street2004ArabicLogicHHL}
\end{proofboundary}
